\Askhsh{25766_SOLUTION}{1}
\begin{ylh}{1}
    \item [Ύλη από εικόνα]
\end{ylh}
ΛΥΣΗ
\begin{alist}
    \item Θεωρούμε τα διαστήματα $\Delta_1 = (-\infty, -2]$, $\Delta_2 = [-2, 0]$, $\Delta_3 = [0, 2]$, $\Delta_4 = [2, +\infty)$. Η συνάρτηση είναι συνεχής σε καθένα από αυτά και από το πρόσημο της παραγώγου της συμπεραίνουμε ότι η $f$ είναι:
    \begin{itemize}
        \item γνησίως αύξουσα σε καθένα από τα διαστήματα $\Delta_1, \Delta_3$
        \item γνησίως φθίνουσα σε καθένα από τα διαστήματα $\Delta_2, \Delta_4$
    \end{itemize}
    Σε ότι αφορά στα ακρότατα της $f$, ισχύει:
    \begin{itemize}
        \item Η συνάρτηση παρουσιάζει (ολικό) μέγιστο για $x=-2, x=2$ και είναι ίσο με $f(-2)=f(2)=5$, αφού η $f$ είναι άρτια.
        \item Η συνάρτηση παρουσιάζει (τοπικό) ελάχιστο για $x=0$, το $f(0)=1$.
    \end{itemize}

    \item Ισχύουν:
    \[ \lim_{x \to +\infty} f(x) = -\infty \]
    και αν θέσουμε $u=-x$,
    \[ \lim_{x \to -\infty} f(x) \overset{f \text{ άρτια}}{=} \lim_{x \to -\infty} f(-x) \overset{u=-x}{=} \lim_{u \to +\infty} f(u) = -\infty \]
    Έτσι, λόγω της συνέχειας της $f$ και σε συνδυασμό με την μονοτονία της, έχουμε:
    \[ f(\Delta_1) = \left(\lim_{x \to -\infty} f(x), f(-2)\right] = (-\infty, 5] \]
    \[ f(\Delta_2) = [f(0), f(-2)] = [1, 5] \]
    \[ f(\Delta_3) = [f(0), f(2)] = [1, 5] \]
    \[ f(\Delta_4) = \left(\lim_{x \to +\infty} f(x), f(2)\right] = (-\infty, 5] \]
    Επομένως, το σύνολο τιμών της $f$ είναι:
    \[ f(\mathbb{R}) = f(\Delta_1 \cup \Delta_2 \cup \Delta_3 \cup \Delta_4) = f(\Delta_1) \cup f(\Delta_2) \cup f(\Delta_3) \cup f(\Delta_4) = (-\infty, 5]. \]

    \item Από το ερώτημα β) προκύπτει ότι $f(x) \le 5$ με την ισότητα να ισχύει μόνο για $x=-2$ και $x=2$. Επιπλέον, αν θεωρήσουμε $g(x)=|x^2-4|+5$, τότε $g(x) \ge 5$ με την ισότητα να ισχύει μόνο για $x=-2$ και $x=2$.
    Επομένως οι αριθμοί $-2, 2$ είναι λύσεις της εξίσωσης $f(x)=g(x)$ και είναι οι μοναδικές, αφού για οποιαδήποτε άλλη τιμή του $x$ ισχύει $f(x)<g(x)$, οπότε η εξίσωση $f(x)=g(x)$ είναι αδύνατη.

    \item Αν θέσουμε $u=-x$, τότε έχουμε $du=-dx$ και τα νέα όρια ολοκλήρωσης είναι $u_1=1$ για $x=-1$ και $u_2=-1$ για $x=1$, οπότε το ολοκλήρωμα $I$ γράφεται:
    \[ I = \int_{-1}^{1} xf(x)dx \overset{f \text{ άρτια}}{=} \int_{-1}^{1} xf(-x)dx \overset{u=-x}{=} \int_{1}^{-1} u f(u)(-du) = -\int_{1}^{-1} uf(u)du = \int_{-1}^{1} uf(u)du \]
    Άρα $I = -I$, οπότε $2I=0$, άρα $I=0$.
\end{alist}